% Figure: the far field of a bounded charge distribution falls off faster
% the higher its lowest non-vanishing multipole. A net charge (monopole)
% falls as 1/r^2, a neutral pair (dipole) as 1/r^3, a quadrupole as 1/r^4.
% The plot is the field magnitude on log-log axes so each power law is a
% straight line of slope equal to the negative exponent.
\begin{figure}[htbp]
\centering
\begin{tikzpicture}[>=Latex]
\begin{loglogaxis}[
  width=11.5cm, height=7cm,
  xlabel={distance from source $r/d$ (source size $d$)},
  ylabel={field magnitude (arbitrary units)},
  xmin=1, xmax=100,
  ymin=1e-8, ymax=1,
  legend pos=south west,
  legend cell align=left,
  font=\scriptsize,
  grid=both,
  grid style={enggray!20},
]
% monopole ~ 1/r^2, dipole ~ 1/r^3, quadrupole ~ 1/r^4, normalised to 1 at r/d=1
\addplot[engblue, line width=1.2pt, domain=1:100, samples=60]
  {x^(-2)};
\addlegendentry{monopole, $E\propto r^{-2}$}
\addplot[engred, line width=1.2pt, domain=1:100, samples=60]
  {x^(-3)};
\addlegendentry{dipole, $E\propto r^{-3}$}
\addplot[engamber, line width=1.2pt, domain=1:100, samples=60]
  {x^(-4)};
\addlegendentry{quadrupole, $E\propto r^{-4}$}
\end{loglogaxis}
\end{tikzpicture}
\caption[Multipole falloff of a bounded charge distribution]{Field
magnitude of a bounded charge distribution against distance, on log-log
axes, for the three lowest multipoles. A distribution with net charge
carries a monopole term that falls as $r^{-2}$; a neutral distribution
with separated positive and negative charge has no monopole and its
leading dipole term falls as $r^{-3}$; a distribution with no net charge
and no net dipole, such as two opposed dipoles, has a quadrupole term
falling as $r^{-4}$. Each higher multipole adds one power of the source
size over distance, so far from any source the lowest surviving multipole
dominates and the field looks simpler the farther one stands. The rule
sets how quickly a bounded source stops mattering: a shielded, balanced
cable pair with no net current or charge presents a dipole or quadrupole
to the outside world and its stray field dies far faster than that of an
unbalanced single conductor.}
\label{fig:vol04ch08:dipole-multipole}
\end{figure}
